\documentclass[11pt]{article}
\usepackage[margin=25mm]{geometry}
\usepackage{fontspec}\setmainfont{Latin Modern Roman}
\usepackage{amsmath,amssymb,mathtools}
\usepackage{unicode-math}\setmathfont{Latin Modern Math}
\usepackage{xurl,hyperref}
\hypersetup{colorlinks=true,urlcolor=blue}
\setcounter{secnumdepth}{0}
\setlength{\parindent}{0pt}\setlength{\parskip}{5pt}
\setlength{\emergencystretch}{3em}
\newcommand{\sref}[2]{\hyperlink{source:#1:#2}{[S#2]}}
\newcommand{\cataloguescope}[1]{\paragraph{Recorded target.}#1}
\begin{document}
\section{Kirby Problem 1.36}\label{statement}

\subsection{Definitions and mathematical statement}

\paragraph{Shared notation.}$\mathbb N=\{1,2,\ldots\}$, and $\mathbb Z,\mathbb Q,\mathbb R,\mathbb C$ denote the integers, rationals, reals and complex numbers. Any other conventions are specified locally.


For an oriented surface $\Sigma$, the relevant skein algebra is the module of framed links in $\Sigma\times[0,1]$ modulo the local $\mathfrak{sl}(2)$, $\mathfrak{sl}(N)$ or HOMFLYPT skein relations, with multiplication by stacking in the interval direction. In the $\mathfrak{sl}(2)$ case the Kauffman bracket relations resolve a crossing as $A$ times one smoothing plus $A^{-1}$ times the other, and remove a trivial circle with factor $-A^2-A^{-2}$. The $\mathfrak{sl}(N)$ and HOMFLYPT variants use their corresponding oriented quantum skein relations. A categorification should assign graded link objects whose Grothendieck classes give the skein algebra, together with a categorical multiplication inducing the stacking product; merely assigning link homologies without a compatible product does not fulfill the target.
\paragraph{Statement recorded in the supplied source.}
Construct such categorifications of the $\mathfrak{sl}(2)$, $\mathfrak{sl}(N)$ and HOMFLYPT skein algebras for surfaces, including a categorical analogue of the skein multiplication. \sref{statement}{1}\sref{statement}{2}
\subsection{Short English statement}
Categorify all three specified surface skein algebras, with the stacking product represented categorically.
\subsection{Sources}
\begin{description}
\item[\hypertarget{source:statement:1}{S1}] ULAM.AI and the UnsolvedMath Contributors, \emph{UnsolvedMath: A Curated Collection of Open Mathematics Problems}, catalogue entry. CC BY 4.0. \url{https://huggingface.co/datasets/ulamai/UnsolvedMath}.
\item[\hypertarget{source:statement:2}{S2}] R. \.{I}nan\c{c} Baykur, Robion C. Kirby and Daniel Ruberman, K3 -- A New Problem List in Low-Dimensional Topology, author preliminary edition (2026), Problem 1.36 and its remarks. \url{https://math.berkeley.edu/sites/default/files/surv-295-ruberman-watermarked-author-pdf.pdf}.
\end{description}
\label{end:statement}
\end{document}
